% BOOK-FILE-ID: f-c679a74cbb29
% 第4章 標準：4つの基本パターン・7つの応用パターンを各5題。
% 主分類・出題意図・別解の対応は docs/chapter4-standard-design-2026-10-04.md。

\subsection{標準}\label{節：標準}
\chapterref{章：漸化式の解法}の4つの基本パターン・7つの応用パターンから,各5題,計55題を出題する.
移項や約分で整理し,差や比を取る量を選ぼう.
別解では手掛かり・計算量・条件を比較できる.
特に\bookterm{ratio}{階比}型と\bookterm{logarithm}{対数}型には共通する考え方がある.
以下,$n$は正の整数とし,必要なときは
\[
  S_n=\sum_{k=1}^{n}a_k
\]
とおく.

\subsubsection{問題}\label{節：標準：問題}
\begin{enumerate}[label=\textbf{問\arabic*.},leftmargin=3.8em,format=\bookitemlink{practice-basic}{problem}{answer}]
  % Q1
  \item $a_1=3$ とし,
    \[
      (2a_n+3)a_{n+1}=3a_n
    \]
    を満たす\bookterm{sequence}{数列}$\{a_n\}$の\bookterm{general}{一般項}を求めよ.
    % S1
  \item $S_n=2a_n-n^2$
    を満たす\bookterm{sequence}{数列}$\{a_n\}$の\bookterm{general}{一般項}を求めよ.
    % R1
  \item $a_1=3,\qquad a_{n+1}=2(n+1)a_n$
    で定められる\bookterm{sequence}{数列}の\bookterm{general}{一般項}を求めよ.
    % M1
  \item $a_1=2,\qquad a_{n+1}=2a_n+n-1$
    で定められる\bookterm{sequence}{数列}$\{a_n\}$の\bookterm{general}{一般項}を求めよ.
    % G1
  \item $a_1=6,\qquad a_{n+1}=-\frac12a_n$
    で定められる\bookterm{sequence}{数列}の\bookterm{general}{一般項}を求めよ.
    % T1
  \item $a_1=2,\,a_2=5$ とし,
    \[
      a_{n+2}+6a_n=5a_{n+1}
    \]
    を満たす\bookterm{sequence}{数列}$\{a_n\}$の\bookterm{general}{一般項}を求めよ.
    % U1
  \item $a_1=3,\qquad b_1=1,$
    \[
      \begin{cases}
        a_{n+1}=2a_n-b_n,\\
        b_{n+1}=-a_n+2b_n
      \end{cases}
    \]
    で定められる2つの\bookterm{sequence}{数列}$\{a_n\},\{b_n\}$の\bookterm{general}{一般項}を求めよ.
    % L1
  \item $a_1=4,\qquad a_{n+1}=\frac{a_n^3}{4}$
    で定められる\bookterm{sequence}{数列}$\{a_n\}$の\bookterm{general}{一般項}を求めよ.
    % C1
  \item $a_1=5,\qquad a_{n+1}=3a_n-4$
    で定められる\bookterm{sequence}{数列}$\{a_n\}$の\bookterm{general}{一般項}を求めよ.
    % D1
  \item $a_1=\frac56,\qquad 3a_{n+1}=3a_n-2$
    で定められる\bookterm{sequence}{数列}の\bookterm{general}{一般項}を求めよ.
    % F1
  \item $a_1=2,\qquad a_{n+1}=a_n+3n-1$
    で定められる\bookterm{sequence}{数列}の\bookterm{general}{一般項}を求めよ.
    % T2
  \item $a_1=3,\,a_2=1$ とし,
    \[
      a_{n+2}+a_{n+1}=6a_n
    \]
    を満たす\bookterm{sequence}{数列}$\{a_n\}$の\bookterm{general}{一般項}を求めよ.
    % F2
  \item $a_1=1,\qquad a_{n+1}=a_n+3n^2+3n+1$
    で定められる\bookterm{sequence}{数列}の\bookterm{general}{一般項}を求めよ.
    % Q2
  \item $a_1=2,\,a_{n+1}=-\frac{2a_n}{3a_n+2}$
    で定められる\bookterm{sequence}{数列}$\{a_n\}$の\bookterm{general}{一般項}を求めよ.
    % C2
  \item $a_1=1,\qquad 2a_{n+1}+a_n=9$
    で定められる\bookterm{sequence}{数列}$\{a_n\}$の\bookterm{general}{一般項}を求めよ.
    % U2
  \item $a_1=3,\qquad b_1=1,$
    \[
      \begin{cases}
        a_{n+1}=4a_n-2b_n,\\
        b_{n+1}=a_n+b_n
      \end{cases}
    \]
    で定められる2つの\bookterm{sequence}{数列}$\{a_n\},\{b_n\}$の\bookterm{general}{一般項}を求めよ.
    % D2
  \item $a_1=-4,\qquad 2a_{n+1}+a_n=3a_n+6$
    で定められる\bookterm{sequence}{数列}の\bookterm{general}{一般項}を求めよ.
    % G2
  \item $a_1=3,\qquad 2a_{n+1}-5a_n=a_{n+1}+a_n$
    で定められる\bookterm{sequence}{数列}の\bookterm{general}{一般項}を求めよ.
    % R2
  \item $a_1=6,\qquad
      (n+1)(n+2)a_{n+1}=n(n+3)a_n$
    で定められる\bookterm{sequence}{数列}の\bookterm{general}{一般項}を求めよ.
    % M2
  \item $a_1=3,\qquad a_{n+1}+a_n=2n^2+2n+1$
    で定められる\bookterm{sequence}{数列}$\{a_n\}$の\bookterm{general}{一般項}を求めよ.
    % S2
  \item $S_n=a_n+(n-1)2^{n-1}$
    を満たす\bookterm{sequence}{数列}$\{a_n\}$の\bookterm{general}{一般項}を求めよ.
    % L2
  \item $a_1=5,\qquad a_{n+1}=a_n^2-4a_n+6$
    で定められる\bookterm{sequence}{数列}$\{a_n\}$の\bookterm{general}{一般項}を求めよ.
    % T3
  \item $a_1=1,\,a_2=0$ とし,
    \[
      a_{n+2}-4a_{n+1}+4a_n=0
    \]
    を満たす\bookterm{sequence}{数列}$\{a_n\}$の\bookterm{general}{一般項}を求めよ.
    % M3
  \item $a_1=1,\qquad a_{n+1}=3a_n-2^n$
    で定められる\bookterm{sequence}{数列}$\{a_n\}$の\bookterm{general}{一般項}を求めよ.
    % U3
  \item $a_1=1,\qquad b_1=2,$
    \[
      \begin{cases}
        a_{n+1}=2a_n+b_n,\\
        b_{n+1}=2b_n
      \end{cases}
    \]
    で定められる2つの\bookterm{sequence}{数列}$\{a_n\},\{b_n\}$の\bookterm{general}{一般項}を求めよ.
    % Q3
  \item $a_1=4,\,a_{n+1}=\frac{3a_n+2}{a_n+2}$
    で定められる\bookterm{sequence}{数列}$\{a_n\}$の\bookterm{general}{一般項}を求めよ.
    % G3
  \item $a_1=4,\qquad (n+1)a_{n+1}=\frac{n+1}{3}a_n$
    で定められる\bookterm{sequence}{数列}の\bookterm{general}{一般項}を求めよ.
    % D3
  \item $a_1=1-\sqrt2,\qquad a_{n+1}+n=a_n+n+\sqrt2$
    で定められる\bookterm{sequence}{数列}の\bookterm{general}{一般項}を求めよ.
    % L3
  \item $a_1=2,\qquad a_{n+1}=a_n^{\frac{n+2}{n}}$
    で定められる正の\bookterm{sequence}{数列}$\{a_n\}$の\bookterm{general}{一般項}を求めよ.
    % S3
  \item $S_1=1,\qquad S_{n+1}=3S_n+2$
    を満たす\bookterm{sequence}{数列}$\{a_n\}$の\bookterm{general}{一般項}を求めよ.
    % F3
  \item $a_1=3,\qquad a_{n+1}=a_n+2\cdot3^n$
    で定められる\bookterm{sequence}{数列}の\bookterm{general}{一般項}を求めよ.
    % R3
  \item $a_1=4,\qquad (n+1)^2a_{n+1}=n^2a_n$
    で定められる\bookterm{sequence}{数列}の\bookterm{general}{一般項}を求めよ.
    % C3
  \item $a_1=0,\qquad a_{n+1}-a_n=\frac{6-a_n}{3}$
    で定められる\bookterm{sequence}{数列}$\{a_n\}$の\bookterm{general}{一般項}を求めよ.
    % S4
  \item $a_1=1,\qquad a_{n+1}=2S_n+2^n$
    で定められる\bookterm{sequence}{数列}$\{a_n\}$の\bookterm{general}{一般項}を求めよ.
    % Q4
  \item $a_1=1,\,a_{n+1}=2+\frac3{a_n}$
    で定められる\bookterm{sequence}{数列}$\{a_n\}$の\bookterm{general}{一般項}を求めよ.
    % G4
  \item \bookterm{sequence}{数列}$\{a_n\}$は,ある実数$r$について
    \[
      a_{n+1}=ra_n
    \]
    を満たし,$a_2=-6$, $a_5=48$である.$r$および\bookterm{general}{一般項}$a_n$を求めよ.
    % F4
  \item $a_1=0,\qquad a_{n+1}=a_n+\frac{2}{(n+1)(n+3)}$
    で定められる\bookterm{sequence}{数列}の\bookterm{general}{一般項}を求めよ.
    % T4
  \item $a_1=3,\,a_2=5$ とし,
    \[
      a_{n+2}=2a_{n+1}+3a_n-8
    \]
    を満たす\bookterm{sequence}{数列}$\{a_n\}$の\bookterm{general}{一般項}を求めよ.
    % U4
  \item $a_1=0,\qquad b_1=1,$
    \[
      \begin{cases}
        a_{n+1}=2a_n+b_n+1,\\
        b_{n+1}=a_n+2b_n-1
      \end{cases}
    \]
    で定められる2つの\bookterm{sequence}{数列}$\{a_n\},\{b_n\}$の\bookterm{general}{一般項}を求めよ.
    % L4
  \item $a_1=1,\qquad a_{n+1}=2\sqrt{a_n}$
    で定められる\bookterm{sequence}{数列}$\{a_n\}$の\bookterm{general}{一般項}を求めよ.
    % C4
  \item 実数$t$に対し,
    \[
      a_1=t,\qquad a_{n+1}=4a_n-9
    \]
    で\bookterm{sequence}{数列}$\{a_n\}$を定める.\bookterm{general}{一般項}を求め,この\bookterm{sequence}{数列}が定数\bookterm{sequence}{数列}となる$t$の値を求めよ.
    % R4
  \item $a_1=2,\qquad a_{n+1}=2^{2n+1}a_n$
    で定められる\bookterm{sequence}{数列}の\bookterm{general}{一般項}を求めよ.
    % D4
  \item \bookterm{sequence}{数列}$\{a_n\}$は,ある実数$d$について
    \[
      a_{n+1}=a_n+d
    \]
    を満たし,$a_3=8$, $a_8=-7$である.$d$, $a_1$および\bookterm{general}{一般項}$a_n$を求めよ.
    % M4
  \item $a_1=2,\qquad a_{n+1}=2a_n+3\cdot2^n$
    で定められる\bookterm{sequence}{数列}$\{a_n\}$の\bookterm{general}{一般項}を求めよ.
    % U5
  \item $a_1=1,\qquad b_1=0,$
    \[
      \begin{cases}
        a_{n+1}=a_n+2b_n,\\
        b_{n+1}=a_n
      \end{cases}
    \]
    で定められる2つの\bookterm{sequence}{数列}$\{a_n\},\{b_n\}$の\bookterm{general}{一般項}を求めよ.
    % R5
  \item $a_1=1,\qquad a_{n+1}=\frac{n(n+1)}2a_n$
    で定められる\bookterm{sequence}{数列}の\bookterm{general}{一般項}を求めよ.
    % Q5
  \item $a_1=2,\,a_{n+1}=2-\frac1{a_n}$
    で定められる\bookterm{sequence}{数列}$\{a_n\}$の\bookterm{general}{一般項}を求めよ.
    % G5
  \item $a_1=0,\qquad a_{n+1}=7a_n$
    で定められる\bookterm{sequence}{数列}の\bookterm{general}{一般項}を求めよ.
    % L5
  \item $a_1=4,\qquad a_{n+1}=\frac{a_n^2}{2^{n-1}}$
    で定められる\bookterm{sequence}{数列}$\{a_n\}$の\bookterm{general}{一般項}を求めよ.
    % F5
  \item $a_1=1,\qquad a_{n+1}=a_n+(-1)^n$
    で定められる\bookterm{sequence}{数列}の\bookterm{general}{一般項}を求めよ.
    % C5
  \item 実数$p,q$を用いて,
    \[
      a_{n+1}=pa_n+q
    \]
    で定められる\bookterm{sequence}{数列}$\{a_n\}$について,
    \[
      a_1=8,\qquad a_2=5,\qquad a_3=\frac72
    \]
    であることが分かっている.$p,q$の値と\bookterm{general}{一般項}を求めよ.
    % T5
  \item $a_1=1,\,a_2=2$ とし,
    \[
      a_{n+2}-3a_{n+1}+2a_n=1
    \]
    を満たす\bookterm{sequence}{数列}$\{a_n\}$の\bookterm{general}{一般項}を求めよ.
    % D5
  \item $a_1=\frac74,\qquad 4a_{n+1}-4a_n=0$
    で定められる\bookterm{sequence}{数列}の\bookterm{general}{一般項}を求めよ.
    % M5
  \item $a_1=0,\qquad a_{n+1}=2a_n+3(-1)^n$
    で定められる\bookterm{sequence}{数列}$\{a_n\}$の\bookterm{general}{一般項}を求めよ.
    % S5
  \item $a_1=2,\qquad S_n=\frac{n+2}{3}a_n$
    を満たす\bookterm{sequence}{数列}$\{a_n\}$の\bookterm{general}{一般項}を求めよ.
\end{enumerate}

\subsubsection{方針}\label{節：標準：方針}
\begin{enumerate}[label=\textbf{問\arabic*.},leftmargin=3.8em,format=\bookitemlink{practice-basic}{hint}{problem}]
  % Q1
  \item $a_{n+1}$について解き,両辺の逆数を取る.$\frac{1}{a_n}$の増え方に注目する.
    % S1
  \item $n=1$から初項を求める. 次に$S_{n+1}-S_n=a_{n+1}$を使い, 総和を消す. 得られた\bookterm{recurrence}{漸化式}では$a_n$に$n$の一次式を加えて\bookterm{geometric}{等比数列}を作る.
    % R1
  \item $(n+1)!=(n+1)n!$を使い,$a_n$を$n!$で割る.
    % M1
  \item $a_n$から適切な一次式を引き,残りが\bookterm{geometric}{等比数列}になるようにする.
    % G1
  \item 初項から第$n$項までに,$-\frac12$を掛ける回数を考える.
    % T1
  \item \bookterm{characteristic}{特性方程式}を解き,$a_{n+1}-2a_n$と$a_{n+1}-3a_n$を求める.二式の差で$a_n$が残る.
    % U1
  \item 2つの式を足すと和は変わらず, 引くと差は3倍になる. 和と差を求めてから$a_n,b_n$に戻す.
    % L1
  \item 全項が正であることを確かめ,底を$2$とする\bookterm{logarithm}{対数}を取る.
    % C1
  \item $a_n$の値が変わらないとしたときの値を求め,その値との差を考える.
    % D1
  \item 両辺を$3$で割り,$a_{n+1}-a_n$を求める.
    % F1
  \item $a_{k+1}-a_k=3k-1$を$k=1$から$n-1$まで加える.
    % T2
  \item \bookterm{characteristic}{特性方程式}の解は$2$と$-3$である.負の解に対応する\bookterm{auxiliary}{補助数列}と$(-3)^{n-1}$の符号に注意する.
    % F2
  \item $3n^2+3n+1$が$(n+1)^3-n^3$であることを使う.
    % Q2
  \item 逆数を取って一次の\bookterm{recurrence}{漸化式}に帰着させる.最初の数項を計算すると,より短い解法も見つかる.
    % C2
  \item まず$a_{n+1}$について解く.\bookterm{fixed}{不動点}との差は負の公比をもつ\bookterm{geometric}{等比数列}になる.
    % U2
  \item $a_n+t b_n$が\bookterm{geometric}{等比数列}になるように$t$を選ぶ. $t=-1,-2$で2本の\bookterm{geometric}{等比数列}を得られる. 片方を消して隣接3項間の\bookterm{recurrence}{漸化式}にする別解もある.
    % D2
  \item 両辺から$a_n$を引き,その後で$2$で割る.
    % G2
  \item $a_{n+1}$を左辺に,$a_n$を右辺にまとめる.
    % R2
  \item 係数を$\frac{n}{n+1}\frac{n+3}{n+2}$と分解し,順に掛ける.
    % M2
  \item 右辺を$(n+1)^2+n^2$と見れば,引くべき二次式が分かる.
    % S2
  \item $S_{n+1}-S_n=a_{n+1}$に代入する. 両辺の$a_{n+1}$が消え,$a_n$が直接求まる. $n=1$の与式だけでは初項が決まらないことにも注意する.
    % L2
  \item \bookterm{square}{平方完成}すると$a_{n+1}-2=(a_n-2)^2$となる.正の量$a_n-2$の\bookterm{logarithm}{対数}を取る.
    % T3
  \item \bookterm{characteristic}{特性方程式}は重解$2$を持つ.$a_{n+1}-2a_n$を\bookterm{geometric}{等比数列}にした後,$\frac{a_n}{2^{n-1}}$を考える.
    % M3
  \item $2^n$に関係する\bookterm{geometric}{等比数列}を引くか,$2^{n-1}$で割って一次の\bookterm{recurrence}{漸化式}にする.
    % U3
  \item $b_n$は$a_n$と関係なく先に解ける. その結果を第1式に代入し,$\frac{a_n}{2^{n-1}}$の\bookterm{difference}{階差}を調べる.
    % Q3
  \item \bookterm{fixed}{不動点}は$2$と$-1$である.$\frac{a_n-2}{a_n+1}$を考える.$\frac{1}{a_n-2}$を考える別解もある.
    % G3
  \item $n$は正整数なので,$n+1$で割ることができる.
    % D3
  \item 両辺に現れる$n$を消去してから,隣り合う項の差を見る.
    % L3
  \item \bookterm{logarithm}{対数}を取った\bookterm{sequence}{数列}の\bookterm{ratio}{階比}が$\frac{n+2}{n}$となる.積の途中の因子が消えることに注目する.
    % S3
  \item $S_n+1$を\bookterm{geometric}{等比数列}として解いてから,$a_n=S_n-S_{n-1}$を使う. この差の公式を使う範囲と$a_1$を区別する.
    % F3
  \item \bookterm{difference}{階差}を加えた後,$3+3^2+\cdots+3^{n-1}$を求める.
    % R3
  \item 左辺は第$n+1$項に$(n+1)^2$を掛けた形,右辺は第$n$項に$n^2$を掛けた形である.
    % C3
  \item 右辺が$0$になる値$6$との差を取ると,等比型になる.
    % S4
  \item $S_{n+1}=S_n+a_{n+1}$から$S_n$の\bookterm{recurrence}{漸化式}を作る. $S_n+2^n$が\bookterm{geometric}{等比数列}になる. 1つ前の式との差から$a_n$の\bookterm{recurrence}{漸化式}を作る別解もある.
    % Q4
  \item 右辺を一つの分数にまとめる.\bookterm{fixed}{不動点}$3,-1$からの比を取ると,公比が負の\bookterm{geometric}{等比数列}になる.
    % G4
  \item 第$2$項から第$5$項までには$r$を$3$回掛ける.
    % F4
  \item $\frac{2}{(k+1)(k+3)}=\frac1{k+1}-\frac1{k+3}$と分解する.
    % T4
  \item 定数解$2$を引く.その後の\bookterm{characteristic}{特性方程式}の一つの解$3$を使うと,初項が$0$の\bookterm{auxiliary}{補助数列}が現れる.
    % U4
  \item 定数項にも注意して2つの式を足し引きする. 和は3倍され, 差には2が加えられる.
    % L4
  \item 平方根を$a_n^{\frac{1}{2}}$と書き,底を$2$とする\bookterm{logarithm}{対数}を取る.
    % C4
  \item \bookterm{fixed}{不動点}との差を取る.その差の初項が$0$となる場合も,分けて確認する.
    % R4
  \item $2$の累乗を掛けると指数が加わる.\bookterm{logarithm}{対数}を取って和に直してもよい.
    % D4
  \item 第$3$項から第$8$項までには,公差を何回加えるかを数える.
    % M4
  \item $2^{n-1}$で割った\bookterm{sequence}{数列}を考える.変形後は等差型になる.
    % U5
  \item 第1式の添字を1つ進め, 第2式を代入して$b_n$を消す. $a_2$はもとの連立から求める. 最後に$b_n$を求めるときは初項も確認する.
    % R5
  \item $(n-1)!n!$の次の項との比が$n(n+1)$であることを使う.
    % Q5
  \item \bookterm{fixed}{不動点}の方程式は重解$1$を持つ.二つの差の比の代わりに,$\frac{1}{a_n-1}$を考える.
    % G5
  \item $0$に$7$を掛けても$0$のままである.項で割らずに考える.
    % L5
  \item \bookterm{logarithm}{対数}を取ると$b_{n+1}=2b_n-n+1$となる.この式を満たす一次式を探す.
    % F5
  \item $(-1)^k$の和は\bookterm{geometric}{等比数列}の和として求められる.最初の数項も調べる.
    % C5
  \item $n=1,2$の漸化式を引けば$q$が消える.係数を求めた後,\bookterm{fixed}{不動点}との差を取る.
    % T5
  \item 定数を引くだけでは右辺を$0$にできない.$d_n=a_{n+1}-a_n$を考え,一次の\bookterm{recurrence}{漸化式}に帰着させる.
    % D5
  \item $a_{n+1}=a_n$なら,項の値は変わらない.
    % M5
  \item $(-1)^n$も指数の形である.$(-1)^{n-1}$で割った\bookterm{sequence}{数列}を考える.
    % S5
  \item $S_{n+1}-S_n=a_{n+1}$から$\frac{a_{n+1}}{a_n}$を求める. 比を順に掛けると, 分子と分母の大部分が打ち消し合う.
\end{enumerate}

\subsubsection{答え}\label{節：標準：答え}
\begin{enumerate}[label=\textbf{問\arabic*.},leftmargin=3.8em,format=\bookitemlink{practice-basic}{answer}{problem}]
  % Q1
  \item \[
      a_n=\frac{3}{2n-1}
    \]
    % S1
  \item \[
      a_n=6\cdot2^{n-1}-2n-3
    \]
    % R1
  \item \[
      a_n=3\cdot2^{n-1}n!
    \]
    % M1
  \item \[
      a_n=3\cdot2^{n-1}-n
    \]
    % G1
  \item \[
      a_n=6\left(-\frac12\right)^{n-1}
    \]
    % T1
  \item \[
      a_n=2^{n-1}+3^{n-1}
    \]
    % U1
  \item \[
      a_n=2+3^{n-1},\qquad b_n=2-3^{n-1}
    \]
    % L1
  \item \[
      a_n=2^{\,1+3^{n-1}}
    \]
    % C1
  \item \[
      a_n=2+3^n
    \]
    % D1
  \item \[
      a_n=\frac56-\frac23(n-1)=\frac{9-4n}{6}
    \]
    % F1
  \item \[
      a_n=\frac{3n^2-5n+6}{2}
    \]
    % T2
  \item \[
      a_n=2^n+(-3)^{n-1}
    \]
    % F2
  \item \[
      a_n=n^3
    \]
    % Q2
  \item \[
      a_n=\frac4{5(-1)^{n-1}-3}
    \]
    % C2
  \item \[
      a_n=3-2\left(-\frac12\right)^{n-1}
    \]
    % U2
  \item \[
      a_n=4\cdot3^{n-1}-2^{n-1},\qquad b_n=2\cdot3^{n-1}-2^{n-1}
    \]
    % D2
  \item \[
      a_n=3n-7
    \]
    % G2
  \item \[
      a_n=3\cdot6^{n-1}
    \]
    % R2
  \item \[
      a_n=\frac{2(n+2)}n
    \]
    % M2
  \item \[
      a_n=n^2+2(-1)^{n-1}
    \]
    % S2
  \item \[
      a_n=(n+1)2^{n-1}
    \]
    % L2
  \item \[
      a_n=2+3^{\,2^{n-1}}
    \]
    % T3
  \item \[
      a_n=(2-n)2^{n-1}
    \]
    % M3
  \item \[
      a_n=2^n-3^{n-1}
    \]
    % U3
  \item \[
      a_n=n2^{n-1},\qquad b_n=2^n
    \]
    % Q3
  \item \[
      a_n=\frac{10\cdot4^{n-1}+2}{5\cdot4^{n-1}-2}
    \]
    % G3
  \item \[
      a_n=4\left(\frac13\right)^{n-1}
    \]
    % D3
  \item \[
      a_n=1+(n-2)\sqrt2
    \]
    % L3
  \item \[
      a_n=2^{\frac{n(n+1)}2}
    \]
    % S3
  \item \[
      a_1=1,\qquad a_n=4\cdot3^{n-2}\quad(n\ge2)
    \]
    % F3
  \item \[
      a_n=3^n
    \]
    % R3
  \item \[
      a_n=\frac4{n^2}
    \]
    % C3
  \item \[
      a_n=6\left\{1-\left(\frac23\right)^{n-1}\right\}
    \]
    % S4
  \item \[
      a_n=2\cdot3^{n-1}-2^{n-1}
    \]
    % Q4
  \item \[
      a_n=\frac{3^n+(-1)^n}{3^{n-1}-(-1)^n}
    \]
    % G4
  \item \[
      r=-2,\qquad a_n=3(-2)^{n-1}
    \]
    % F4
  \item \[
      a_n=\frac56-\frac1{n+1}-\frac1{n+2}
    \]
    % T4
  \item \[
      a_n=2+3^{n-1}
    \]
    % U4
  \item \[
      a_n=\frac{3^{n-1}+2n-3}{2},\qquad b_n=\frac{3^{n-1}-2n+3}{2}
    \]
    % L4
  \item \[
      a_n=2^{\,2-2^{2-n}}
    \]
    % C4
  \item \BookMathLayout{\[
        a_n=3+(t-3)4^{n-1},\qquad
        \text{定数数列となるのは }t=3
      \]}{
      \[
        \begin{aligned}
          a_n&=3+(t-3)4^{n-1},\\
          &\text{定数数列となるのは }t=3
        \end{aligned}
      \]
    }
    % R4
  \item \[
      a_n=2^{n^2}
    \]
    % D4
  \item \[
      d=-3,\qquad a_1=14,\qquad a_n=17-3n
    \]
    % M4
  \item \[
      a_n=(3n-1)2^{n-1}
    \]
    % U5
  \item \BookMathLayout{\[
        a_n=\frac{2^n+(-1)^{n-1}}{3},\qquad b_n=\frac{2^{n-1}-(-1)^{n-1}}{3}
      \]}{
      \[
        \begin{aligned}
          a_n&=\frac{2^n+(-1)^{n-1}}3,\\
          b_n&=\frac{2^{n-1}-(-1)^{n-1}}3
        \end{aligned}
      \]
    }
    % R5
  \item \[
      a_n=\frac{(n-1)!\,n!}{2^{n-1}}
    \]
    % Q5
  \item \[
      a_n=\frac{n+1}{n}
    \]
    % G5
  \item \[
      a_n=0
    \]
    % L5
  \item \[
      a_n=2^{\,n+2^{n-1}}
    \]
    % F5
  \item \[
      a_n=\frac{1+(-1)^{n-1}}2
    \]
    % C5
  \item \[
      p=\frac12,\qquad q=1,\qquad
      a_n=2+6\left(\frac12\right)^{n-1}
    \]
    % T5
  \item \[
      a_n=2^n-n
    \]
    % D5
  \item \[
      a_n=\frac74
    \]
    % M5
  \item \[
      a_n=(-1)^{n-1}-2^{n-1}
    \]
    % S5
  \item \[
      a_n=n(n+1)
    \]
\end{enumerate}

\subsubsection{解説}\label{節：標準：解説}
\begin{enumerate}[label=\textbf{問\arabic*.},leftmargin=3.8em,format=\bookitemlink{practice-basic}{solution}{problem}]
  % Q1
  \item \booknavref{節：逆数型}{〔復習：逆数型〕}\BookMathLayout{}{\\}
    $a_1>0$であり,$a_n>0$なら$a_{n+1}=\frac{3a_n}{2a_n+3}>0$である.
    全項が正なので逆数を取れる.$b_n=\frac{1}{a_n}$とおくと,
    \[
      b_{n+1}=b_n+\frac23,\qquad b_1=\frac13.
    \]
    よって$b_n=\frac{2n-1}{3}$, $a_n=\frac{3}{2n-1}$となる.復元の分母も正である.
    \par\smallskip\noindent\textbf{【別解】}\par
    最初の三項$\frac{3}{1},\frac{3}{3},\frac{3}{5}$から同じ式を予想する.
    初項と一致し,$a_k=\frac{3}{2k-1}$を漸化式へ入れると$a_{k+1}=\frac{3}{2k+1}$となるので,
    \bookterm{induction}{数学的帰納法}で確かめられる.

    % S1
  \item \booknavref{節：総和型}{〔復習：総和型〕}\BookMathLayout{}{\\}
    $n=1$から$a_1=1$となる.総和を消すため隣り合う式を引くと,
    \[
      a_{n+1}=2a_{n+1}-(n+1)^2-2a_n+n^2
    \]
    より$a_{n+1}=2a_n+2n+1$である.
    $n$の項を消すため$b_n=a_n+un+v$とおけば,
    \[
      b_{n+1}-2b_n=(2-u)n+1+u-v.
    \]
    $u=2,v=3$を選ぶと$b_{n+1}=2b_n$, $b_1=6$となる.
    よって$b_n=6\cdot2^{n-1}$から$a_n=6\cdot2^{n-1}-2n-3$を得る.

    % R1
  \item \booknavref{節：階比型}{〔復習：階比型〕}\BookMathLayout{}{\\}
    $(n+1)!=(n+1)n!$に合わせて$b_n=\frac{a_n}{n!}$とおくと,
    \[
      b_{n+1}=2b_n,\quad b_1=3,\quad b_n=3\cdot2^{n-1}.
    \]
    したがって$a_n=3\cdot2^{n-1}n!$である.
    係数$n+1$を約分するため,$(n-1)!$ではなく$n!$で割った.
    \par\smallskip\noindent\textbf{【別解】}\par
    係数を順に掛けても,
    \[
      a_n=3\prod_{k=1}^{n-1}2(k+1)=3\cdot2^{n-1}n!
    \]
    となる.$2$は$n-1$個あり,$n=1$の空積は$1$とする.

    % M1
  \item \booknavref{節：階差等比複合型}{〔復習：階差等比複合型〕}\BookMathLayout{}{\\}
    定数だけでは$n-1$を消せないので,元の漸化式を満たす一次式$g(n)=un+v$を探す.
    $g(n+1)=2g(n)+n-1$の係数を比べると,
    \[
      u=2u+1,\quad u+v=2v-1
    \]
    より$g(n)=-n$である.$b_n=a_n+n$とおくと,
    \[
      b_{n+1}=2b_n,\quad b_1=3,\quad a_n=3\cdot2^{n-1}-n.
    \]
    \par\smallskip\noindent\textbf{【別解：\bookterm{difference}{階差}を先に求める】}\par
    隣り合う2つの漸化式を引くと,$d_n=a_{n+1}-a_n$は
    \[
      d_{n+1}=2d_n+1,\quad d_1=2
    \]
    を満たす.よって$d_n+1=3\cdot2^{n-1}$であり,
    \[
      a_n=2+\sum_{k=1}^{n-1}(3\cdot2^{k-1}-1)
      =3\cdot2^{n-1}-n.
    \]
    空の和を$0$とすれば$n=1$も含む.
    主解は一次式を引いて付加項を消し,別解は差を取って付加項を定数にする.
    主解の方が短いが,一次式の候補が見えないときにも\bookterm{difference}{階差}を試せる.

    % G1
  \item \booknavref{節：等比型}{〔復習：等比型〕}\BookMathLayout{}{\\}
    初項$6$,公比$-\frac{1}{2}$の\bookterm{geometric}{等比数列}だから,
    \[
      a_n=6\left(-\frac12\right)^{n-1}.
    \]
    符号は交互に変わり,絶対値は毎回半分になる.

    % T1
  \item \booknavref{節：隣接3項型}{〔復習：隣接3項型〕}\BookMathLayout{}{\\}
    \bookterm{characteristic}{特性方程式}は$(x-2)(x-3)=0$である.
    $b_n=a_{n+1}-2a_n$, $c_n=a_{n+1}-3a_n$とおくと,
    \[
      b_{n+1}=3b_n,\quad c_{n+1}=2c_n,
    \]
    \[
      b_1=1,\quad c_1=-1
    \]
    より$b_n=3^{n-1}$, $c_n=-2^{n-1}$となる.
    二式を引いて共通の$a_{n+1}$を消せば,
    \[
      a_n=b_n-c_n=2^{n-1}+3^{n-1}.
    \]

    % U1
  \item \booknavref{節：連立型}{〔復習：連立型〕}\BookMathLayout{}{\\}
    $s_n=a_n+b_n$, $d_n=a_n-b_n$とおく.二式を足し引きすると,
    \[
      s_{n+1}=s_n,\quad d_{n+1}=3d_n,\quad s_1=4,\ d_1=2.
    \]
    よって$s_n=4$, $d_n=2\cdot3^{n-1}$であり,
    \[
      a_n=\frac{s_n+d_n}{2}=2+3^{n-1},
    \]
    \[
      b_n=\frac{s_n-d_n}{2}=2-3^{n-1}.
    \]
    \par\smallskip\noindent\textbf{【別解】}\par
    和が一定なので$b_n=4-a_n$を代入すると,$a_{n+1}=3a_n-4$となる.
    $a_{n+1}-2=3(a_n-2)$と$a_1-2=1$から$a_n$を求め,$b_n=4-a_n$で戻す.

    % L1
  \item \booknavref{節：対数型}{〔復習：対数型〕}\BookMathLayout{}{\\}
    $a_1>0$であり,正の項から次も正になるので全項の\bookterm{logarithm}{対数}を取れる.
    指数$3$を係数へ変えるため$b_n=\log_2a_n$とおくと,
    \[
      b_{n+1}=3b_n-2,\quad b_1=2.
    \]
    \bookterm{fixed}{不動点}$1$を引けば$b_{n+1}-1=3(b_n-1)$だから,
    \[
      b_n=1+3^{n-1},\quad a_n=2^{\,1+3^{n-1}}.
    \]
    \par\smallskip\noindent\textbf{【別解：定数倍して置き換える】}\par
    $a_n=Kc_n$なら$c_{n+1}=\left(\frac{K^2}{4}\right)c_n^3$である.
    $c_{n+1}=c_n^3$の形にしたいので,残る係数が$1$になるように$K$を選ぶ.
    $K=2$で係数を消すと$c_{n+1}=c_n^3$, $c_1=2$となるので,
    \[
      c_n=2^{3^{n-1}},\qquad a_n=2c_n.
    \]
    この式は初項に合い,三乗すると次項の式になる.
    \par\smallskip\noindent\textbf{【別解：\bookterm{ratio}{階比}を置く】}\par
    全項が正なので$r_n=\frac{a_{n+1}}{a_n}=\frac{a_n^2}{4}$とおくと,
    \[
      r_{n+1}=r_n^3,\quad r_1=4,\quad r_n=4^{3^{n-1}}.
    \]
    比を掛け戻せば,
    \[
      \begin{aligned}
        a_n&=4\prod_{k=1}^{n-1}r_k
          =4^{\,1+\sum_{k=1}^{n-1}3^{k-1}}\\
        &=2^{\,1+3^{n-1}}.
      \end{aligned}
    \]
    $\frac{a_n}{2}$と置く方法が最短である.
    \bookterm{logarithm}{対数}を取れば,その置換が見えなくても一次の漸化式にできる.
    \bookterm{ratio}{階比}には積を戻す手間があるが,積と指数の和のつながりが見える.

    % C1
  \item \booknavref{節：特性方程式型}{〔復習：特性方程式型〕}\BookMathLayout{}{\\}
    定数項を消すため,$\alpha=3\alpha-4$を解いて$\alpha=2$を引く.
    $b_n=a_n-2$とおくと,
    \[
      b_{n+1}=3b_n,\quad b_1=3,\quad a_n=2+3^n.
    \]
    \par\smallskip\noindent\textbf{【別解：\bookterm{difference}{階差}を使う】}\par
    添字を1つずらした式と引くと,$d_n=a_{n+1}-a_n$は$d_{n+1}=3d_n$, $d_1=6$を満たす.
    よって,
    \[
      a_n=5+6\sum_{k=1}^{n-1}3^{k-1}=2+3^n.
    \]
    空の和を$0$とすれば初項も含む.

    % D1
  \item \booknavref{節：等差型}{〔復習：等差型〕}\BookMathLayout{}{\\}
    両辺を$3$で割ると$a_{n+1}=a_n-\frac{2}{3}$となる.
    公差$-\frac{2}{3}$の\bookterm{arithmetic}{等差数列}なので,
    \[
      a_n=\frac56-\frac23(n-1)=\frac{9-4n}{6}.
    \]

    % F1
  \item \booknavref{節：階差型}{〔復習：階差型〕}\BookMathLayout{}{\\}
    初項に$n-1$回分の増加を加えると,
    \BookMathLayout{\[
        \begin{aligned}
          a_n&=2+\sum_{k=1}^{n-1}(3k-1)\\
          &=2+\frac{3n(n-1)}2-(n-1)
          =\frac{3n^2-5n+6}{2}.
        \end{aligned}
      \]}{\[
        \begin{aligned}
          a_n&=2+\sum_{k=1}^{n-1}(3k-1)\\
          &=2+\frac{3n(n-1)}2-(n-1)\\
          &=\frac{3n^2-5n+6}{2}.
        \end{aligned}
      \]}
    $n=1$の空の和は$0$とする.和の上限を$n$にすると第$n+1$項になる.

    % T2
  \item \booknavref{節：隣接3項型}{〔復習：隣接3項型〕}\BookMathLayout{}{\\}
    \bookterm{characteristic}{特性方程式}は$(x-2)(x+3)=0$である.
    $b_n=a_{n+1}-2a_n$, $c_n=a_{n+1}+3a_n$とおけば,
    \[
      b_{n+1}=-3b_n,\quad c_{n+1}=2c_n,
    \]
    \[
      b_1=-5,\quad c_1=10.
    \]
    したがって$b_n=-5(-3)^{n-1}$, $c_n=10\cdot2^{n-1}$であり,
    \[
      a_n=\frac{c_n-b_n}{5}=2^n+(-3)^{n-1}.
    \]
    負の公比もそのまま\bookterm{geometric}{等比数列}として扱える.

    % F2
  \item \booknavref{節：階差型}{〔復習：階差型〕}\BookMathLayout{}{\\}
    $3n^2+3n+1=(n+1)^3-n^3$なので,
    \[
      a_n=1+\sum_{k=1}^{n-1}\{(k+1)^3-k^3\}=n^3.
    \]
    途中の立方数が相殺するため,平方和の計算を省ける.空の和なら初項も含む.
    \par\smallskip\noindent\textbf{【別解1】}\par
    最初の三項$1,8,27$から$a_n=n^3$を予想する.
    初項に合い,$a_{n+1}=n^3+3n^2+3n+1=(n+1)^3$となるので,
    \bookterm{induction}{数学的帰納法}で示せる.
    \par\smallskip\noindent\textbf{【別解2】}\par
    $a_{n+1}-(n+1)^3=a_n-n^3$と書けば,$a_n-n^3$は初項での値$0$に等しい.

    % Q2
  \item \booknavref{節：逆数型}{〔復習：逆数型〕}\BookMathLayout{}{\\}
    初めの項は$2,-\frac{1}{2},2$となり,前の二値は互いに移り合う.
    その分母$3a_n+2$は$8$と$\frac{1}{2}$であるため,全項が定まり,$0$にもならない.
    $b_n=\frac{1}{a_n}$とおくと,
    \[
      b_{n+1}=-b_n-\frac32,\quad b_1=\frac12.
    \]
    $\frac{3}{4}$を加えて等比型にすれば,
    \[
      b_n+\frac34=\frac54(-1)^{n-1},\quad
      a_n=\frac4{5(-1)^{n-1}-3}.
    \]
    復元の分母は奇数番目で$2$,偶数番目で$-8$となる.
    \par\smallskip\noindent\textbf{【別解】}\par
    二値の往復だけで,
    \[
      a_n=\begin{cases}
        2&(n\text{が奇数}),\\
        -\frac{1}{2}&(n\text{が偶数})
      \end{cases}
    \]
    と表せる.逆数の解法は,これを$(-1)^{n-1}$で一式にまとめたものになる.

    % C2
  \item \booknavref{節：特性方程式型}{〔復習：特性方程式型〕}\BookMathLayout{}{\\}
    整理すると$a_{n+1}=-\frac{a_n}{2}+\frac{9}{2}$である.
    \bookterm{fixed}{不動点}$3$を引いて$b_n=a_n-3$とおくと,
    \[
      b_{n+1}=-\frac12b_n,\quad b_1=-2,
    \]
    より$a_n=3-2\left(-\frac{1}{2}\right)^{n-1}$となる.負の公比により$3$との差の符号が交互に変わる.

    % U2
  \item \booknavref{節：連立型}{〔復習：連立型〕}\BookMathLayout{}{\\}
    $a_n+t b_n$を\bookterm{geometric}{等比数列}にするため,
    \[
      a_{n+1}+t b_{n+1}=(4+t)a_n+(t-2)b_n
    \]
    を$(4+t)(a_n+t b_n)$に合わせる.
    $t-2=t(4+t)$より$t=-1,-2$であり,初項を使うと,
    \[
      a_n-b_n=2\cdot3^{n-1},\quad a_n-2b_n=2^{n-1}.
    \]
    二式を解けば,
    \[
      \begin{aligned}
        a_n&=4\cdot3^{n-1}-2^{n-1},\\
        b_n&=2\cdot3^{n-1}-2^{n-1}.
      \end{aligned}
    \]
    \par\smallskip\noindent\textbf{【別解】}\par
    $2b_n=4a_n-a_{n+1}$を使って消去すると,
    \[
      a_{n+2}=4a_{n+1}-2(a_n+b_n)=5a_{n+1}-6a_n.
    \]
    \bookterm{characteristic}{特性方程式}の根$2,3$から$a_n=C2^{n-1}+D3^{n-1}$とおく.
    $a_1=3,a_2=10$より$C=-1,D=4$となり,$b_n=\frac{4a_n-a_{n+1}}{2}$で戻せる.

    % D2
  \item \booknavref{節：等差型}{〔復習：等差型〕}\BookMathLayout{}{\\}
    整理すると$a_{n+1}=a_n+3$である.
    初項$-4$,公差$3$の\bookterm{arithmetic}{等差数列}だから$a_n=-4+3(n-1)=3n-7$となる.

    % G2
  \item \booknavref{節：等比型}{〔復習：等比型〕}\BookMathLayout{}{\\}
    同じ項をまとめると$a_{n+1}=6a_n$となる.
    初項$3$,公比$6$の\bookterm{geometric}{等比数列}だから$a_n=3\cdot6^{n-1}$である.

    % R2
  \item \booknavref{節：階比型}{〔復習：階比型〕}\BookMathLayout{}{\\}
    $n\ge1$なので係数の分母は0でない.係数を2つに分けて掛けると,
    \[
      \begin{aligned}
        a_n&=6\prod_{k=1}^{n-1}\frac{k}{k+1}
        \prod_{k=1}^{n-1}\frac{k+3}{k+2}\\
        &=6\cdot\frac1n\cdot\frac{n+2}3
        =\frac{2(n+2)}n.
      \end{aligned}
    \]
    中間の因子が相殺し,両端だけが残る.$n=1$でも初項に合う.
    \par\smallskip\noindent\textbf{【別解】}\par
    相殺後の因子を組み込み,$b_n=\frac{na_n}{n+2}$とおくと,
    \[
      b_{n+1}=\frac{n+1}{n+3}
      \frac{n(n+3)}{(n+1)(n+2)}a_n=b_n.
    \]
    $b_1=2$なので,$a_n=\frac{2(n+2)}{n}$となる.

    % M2
  \item \booknavref{節：階差等比複合型}{〔復習：階差等比複合型〕}\BookMathLayout{}{\\}
    右辺が$(n+1)^2+n^2$であることから,
    \[
      a_{n+1}-(n+1)^2=-(a_n-n^2).
    \]
    $b_n=a_n-n^2$は$b_{n+1}=-b_n$, $b_1=2$を満たすので,
    \[
      a_n=n^2+2(-1)^{n-1}.
    \]
    二次式を引いた残りが$2,-2$を繰り返す.

    % S2
  \item \booknavref{節：総和型}{〔復習：総和型〕}\BookMathLayout{}{\\}
    $n=1$の式は$S_1=a_1$だけなので初項はまだ決まらない.
    全添字での関係を使って差を取ると,
    \[
      \begin{aligned}
        a_{n+1}&=S_{n+1}-S_n\\
        &=a_{n+1}+n2^n-a_n-(n-1)2^{n-1}.
      \end{aligned}
    \]
    両辺の$a_{n+1}$を消すと$a_n=(n+1)2^{n-1}$となる.
    この計算は$n\ge1$で使え,$a_1=2$も決まる.

    % L2
  \item \booknavref{節：対数型}{〔復習：対数型〕}\BookMathLayout{}{\\}
    和のまま\bookterm{logarithm}{対数}を取る前に\bookterm{square}{平方完成}すると,
    \[
      a_{n+1}-2=(a_n-2)^2.
    \]
    $a_1-2=3>0$から全項で$a_n-2>0$である.
    $b_n=\log_3(a_n-2)$は$b_{n+1}=2b_n$, $b_1=1$を満たすので,
    \[
      b_n=2^{n-1},\qquad a_n=2+3^{\,2^{n-1}}.
    \]
    \par\smallskip\noindent\textbf{【別解：二乗を繰り返す】}\par
    $c_n=a_n-2$なら$c_{n+1}=c_n^2$, $c_1=3$である.
    指数が毎回二倍になるため$c_n=3^{2^{n-1}}$となる.
    初項に合い,二乗して次項になるので\bookterm{induction}{数学的帰納法}でも確かめられる.

    % T3
  \item \booknavref{節：隣接3項型}{〔復習：隣接3項型〕}\BookMathLayout{}{\\}
    \bookterm{characteristic}{特性方程式}は$(x-2)^2=0$である.
    異なる2本の\bookterm{auxiliary}{補助数列}は作れないので,$b_n=a_{n+1}-2a_n$だけを考える.
    \[
      b_{n+1}=2b_n,\quad b_1=-2,\quad b_n=-2^n.
    \]
    $a_{n+1}-2a_n=-2^n$を$2^n$で割り,$c_n=\frac{a_n}{2^{n-1}}$とおくと,
    \[
      c_{n+1}-c_n=-1,\quad c_1=1,
    \]
    より$c_n=2-n$, $a_n=(2-n)2^{n-1}$を得る.
    項$a_2=0$で割っておらず,分母$2^n$は常に0でない.
    \par\smallskip\noindent\textbf{【別解】}\par
    初めから同じ$c_n$を使って元の式を$2^{n+1}$で割ると,
    \[
      c_{n+2}-2c_{n+1}+c_n=0.
    \]
    $c_1=1,c_2=0$なので公差$-1$の\bookterm{arithmetic}{等差数列}となる.
    重解の指数部分を除くと等差型が残ることが直接見える.

    % M3
  \item \booknavref{節：階差等比複合型}{〔復習：階差等比複合型〕}\BookMathLayout{}{\\}
    付加項$-2^n$を消すため,元の漸化式を満たす$g(n)=c2^n$を探す.
    $c2^{n+1}=3c2^n-2^n$から$c=1$である.
    $b_n=a_n-2^n$とおくと$b_{n+1}=3b_n$, $b_1=-1$なので,
    \[
      a_n=2^n-3^{n-1}.
    \]
    \par\smallskip\noindent\textbf{【別解：指数部分で割る】}\par
    $c_n=\frac{a_n}{2^{n-1}}$とおけば$c_{n+1}=\frac{3c_n}{2}-1$, $c_1=1$となる.
    \bookterm{fixed}{不動点}$2$を引くと$c_n-2=-\left(\frac{3}{2}\right)^{n-1}$であり,
    \[
      a_n=2^{n-1}\left\{2-\left(\frac32\right)^{n-1}\right\}.
    \]
    主解は付加項を消し,別解は$2^{n-1}$で割って付加項を定数に変える.
    付加項の底と係数が等しい場合は別解が等差型になる一方,主解には$c2^n$だけでは足りない.
    \sectionref*{節：階差等比複合型}の$p=q$の場合に対応する.

    % U3
  \item \booknavref{節：連立型}{〔復習：連立型〕}\BookMathLayout{}{\\}
    第2式だけで$b_n=2^n$が求まる.これを第1式へ入れ,
    $c_n=\frac{a_n}{2^{n-1}}$とおくと,
    \[
      c_{n+1}=c_n+1,\quad c_1=1.
    \]
    よって$c_n=n$, $a_n=n2^{n-1}$である.
    \par\smallskip\noindent\textbf{【別解】}\par
    $b_n=a_{n+1}-2a_n$を$b_{n+1}=2b_n$へ代入すると,
    \[
      a_{n+2}=4a_{n+1}-4a_n.
    \]
    \bookterm{characteristic}{特性方程式}は$(r-2)^2=0$であり,指数を除けば$c_{n+2}-2c_{n+1}+c_n=0$となる.
    $a_1=1,a_2=4$より$c_1=1,c_2=2$なので,$c_n=n$を得る.

    % Q3
  \item \booknavref{節：一次分数型}{〔復習：一次分数型〕}\BookMathLayout{}{\\}
    \bookterm{fixed}{不動点}は$x^2-x-2=0$の根$2,-1$である.
    $a_1>2$と
    \[
      a_{n+1}-2=\frac{a_n-2}{a_n+2}
    \]
    より全項で$a_n>2$となり,必要な分母は0でない.
    $b_n=\frac{a_n-2}{a_n+1}$とおくと,
    \[
      b_{n+1}=\frac14b_n,\quad b_1=\frac25,
      \quad b_n=\frac25\left(\frac14\right)^{n-1}.
    \]
    $0<b_n\le\frac{2}{5}$に注意して戻せば,
    \[
      a_n=\frac{2+b_n}{1-b_n}
      =\frac{10\cdot4^{n-1}+2}{5\cdot4^{n-1}-2}.
    \]
    \par\smallskip\noindent\textbf{【別解】}\par
    $a_n>2$なので差の逆数$c_n=\frac{1}{a_n-2}$も使える.
    \[
      c_{n+1}=1+4c_n,\quad c_1=\frac12
    \]
    より$c_n+\frac{1}{3}=\left(\frac{5}{6}\right)4^{n-1}$となり,$a_n=2+\frac{1}{c_n}$で戻せる.
    二つの\bookterm{fixed}{不動点}からの差の比と,一つからの差の逆数が,同じ分数型を簡単にしている.

    % G3
  \item \booknavref{節：等比型}{〔復習：等比型〕}\BookMathLayout{}{\\}
    $n\ge1$なので$n+1$で割ると$a_{n+1}=\frac{a_n}{3}$となる.
    見た目の係数に$n$があっても公比は一定であり,$a_n=4\left(\frac{1}{3}\right)^{n-1}$である.

    % D3
  \item \booknavref{節：等差型}{〔復習：等差型〕}\BookMathLayout{}{\\}
    両辺の$n$を消すと$a_{n+1}=a_n+\sqrt2$となる.
    公差$\sqrt2$の\bookterm{arithmetic}{等差数列}だから,
    \[
      a_n=1-\sqrt2+(n-1)\sqrt2=1+(n-2)\sqrt2.
    \]

    % L3
  \item \booknavref{節：対数型}{〔復習：対数型〕}\BookMathLayout{}{\\}
    全項が正なので$b_n=\log_2a_n$とおくと,
    \[
      b_{n+1}=\frac{n+2}{n}b_n,\quad b_1=1.
    \]
    \bookterm{ratio}{階比}型として積を取れば,
    \[
      b_n=\prod_{k=1}^{n-1}\frac{k+2}{k}=\frac{n(n+1)}2,
      \qquad a_n=2^{\frac{n(n+1)}2}.
    \]
    空積を$1$とすれば初項も含む.
    \par\smallskip\noindent\textbf{【別解：累乗して定数にする】}\par
    得られた因子を使って$c_n=a_n^{\dfrac{2}{n(n+1)}}$とおくと,
    \[
      c_{n+1}
      =\left(a_n^{\frac{n+2}{n}}\right)^{\frac2{(n+1)(n+2)}}
      =c_n.
    \]
    $c_1=2$から同じ答えを得る.
    \bookterm{logarithm}{対数}を取ると現れる係数を,別解では累乗の指数で打ち消した.
    どちらも正値が条件となり,別解の指数は主解で割った因子から見つかる.

    % S3
  \item \booknavref{節：総和型}{〔復習：総和型〕}\BookMathLayout{}{\\}
    まず$S_{n+1}+1=3(S_n+1)$, $S_1+1=2$から$S_n=2\cdot3^{n-1}-1$を求める.
    $n\ge2$では,
    \[
      a_n=S_n-S_{n-1}=4\cdot3^{n-2}.
    \]
    初項は$a_1=S_1=1$であり,上の式を$n=1$まで延長できない.
    \par\smallskip\noindent\textbf{【別解】}\par
    $n\ge2$で,総和についての隣り合う2つの式を引けば,
    \[
      a_{n+1}=(3S_n+2)-(3S_{n-1}+2)=3a_n.
    \]
    $S_2=5$より$a_2=4$なので同じ式を得る.
    添字を1つ戻した式を使ったため$n\ge2$であり,$a_1$は別に付け加える.

    % F3
  \item \booknavref{節：階差型}{〔復習：階差型〕}\BookMathLayout{}{\\}
    \bookterm{difference}{階差}を足すと,
    \[
      a_n=3+2\sum_{k=1}^{n-1}3^k=3^n.
    \]
    空の和を$0$とすれば初項も含む.
    \par\smallskip\noindent\textbf{【別解1】}\par
    $2\cdot3^n=3^{n+1}-3^n$から,$a_{n+1}-3^{n+1}=a_n-3^n$となる.
    この差は初項で$0$なので$a_n=3^n$である.
    \par\smallskip\noindent\textbf{【別解2】}\par
    同じ式を予想しても,初項に合い,$a_{n+1}=3^n+2\cdot3^n=3^{n+1}$となるので帰納法で示せる.

    % R3
  \item \booknavref{節：階比型}{〔復習：階比型〕}\BookMathLayout{}{\\}
    両辺が「添字の平方×その項」になっているので,$b_n=n^2a_n$とおく.
    $b_{n+1}=b_n$, $b_1=4$より$a_n=\frac{4}{n^2}$となる.$n\ge1$なので割る因子は0でない.
    \par\smallskip\noindent\textbf{【別解】}\par
    最初の項$4,1,\frac{4}{9},\frac{1}{4}$から同じ式を予想する.
    初項と一致し,$a_{n+1}=\frac{n^2a_n}{(n+1)^2}=\frac{4}{(n+1)^2}$となるので,
    \bookterm{induction}{数学的帰納法}で確かめられる.

    % C3
  \item \booknavref{節：特性方程式型}{〔復習：特性方程式型〕}\BookMathLayout{}{\\}
    右辺に$a_n$が残るため,既知の\bookterm{difference}{階差}として足すことはできない.
    整理すると$a_{n+1}=\frac{2a_n}{3}+2$であり,\bookterm{fixed}{不動点}$6$を引けば,
    \[
      a_{n+1}-6=\frac23(a_n-6).
    \]
    $a_1-6=-6$から$a_n=6-6\left(\frac{2}{3}\right)^{n-1}$となる.
    \par\smallskip\noindent\textbf{【別解：増加量を求める】}\par
    添字を1つずらした式と引くと,$d_n=a_{n+1}-a_n$は$d_{n+1}=\frac{2d_n}{3}$, $d_1=2$を満たす.
    よって,
    \[
      a_n=\sum_{k=1}^{n-1}2\left(\frac23\right)^{k-1}
      =6\left\{1-\left(\frac23\right)^{n-1}\right\}.
    \]
    空の和は$0$とする.

    % S4
  \item \booknavref{節：総和型}{〔復習：総和型〕}\BookMathLayout{}{\\}
    $S_{n+1}=S_n+a_{n+1}=3S_n+2^n$である.
    $c_n=S_n+2^n$なら$c_{n+1}=3c_n$, $c_1=3$となるので$S_n=3^n-2^n$を得る.
    \[
      a_n=S_n-S_{n-1}=2\cdot3^{n-1}-2^{n-1}
      \quad(n\ge2).
    \]
    最終式は$n=1$でも$a_1=1$に合う.
    \par\smallskip\noindent\textbf{【別解】}\par
    $n\ge2$で与式を一つずらして引くと,
    \[
      a_{n+1}-a_n=2(S_n-S_{n-1})+2^{n-1}
    \]
    より$a_{n+1}=3a_n+2^{n-1}$となる.
    $a_2=4=3a_1+1$なので$n=1$でも成り立つ.
    $b_n=a_n+2^{n-1}$は$b_{n+1}=3b_n$, $b_1=2$を満たし,同じ答えを得る.

    % Q4
  \item \booknavref{節：一次分数型}{〔復習：一次分数型〕}\BookMathLayout{}{\\}
    正の項から次も正になるので全項で$a_n>0$であり,元の分母は0でない.
    $a_{n+1}=\frac{2a_n+3}{a_n}$の\bookterm{fixed}{不動点}は$3,-1$である.
    $b_n=\frac{a_n-3}{a_n+1}$とおくと,
    \[
      b_{n+1}=-\frac13b_n,\quad b_1=-1,
      \quad b_n=\frac{(-1)^n}{3^{n-1}}.
    \]
    復元すると,
    \[
      a_n=\frac{3+b_n}{1-b_n}
      =\frac{3^n+(-1)^n}{3^{n-1}-(-1)^n}.
    \]
    最後の分母は奇数番目では$3^{n-1}+1$,偶数番目では$n\ge2$より$3^{n-1}-1>0$となる.
    \par\smallskip\noindent\textbf{【別解】}\par
    $u_1=u_2=1$, $u_{n+2}=2u_{n+1}+3u_n$とすると全項が正であり,
    \[
      \frac{u_{n+2}}{u_{n+1}}=2+\cfrac3{\cfrac{u_{n+1}}{u_n}}.
    \]
    初項の比も$1$なので$a_n=\frac{u_{n+1}}{u_n}$である.
    \bookterm{characteristic}{特性方程式}の根$3,-1$から,
    \BookMathLayout{\[
        u_{n+1}+u_n=2\cdot3^{n-1},
        \quad u_{n+1}-3u_n=-2(-1)^{n-1}.
      \]}{\[
        \begin{aligned}
          u_{n+1}+u_n&=2\cdot3^{n-1},\\
          u_{n+1}-3u_n&=-2(-1)^{n-1}.
        \end{aligned}
      \]}
    二式を引けば$u_n=\frac{3^{n-1}+(-1)^{n-1}}{2}$となり,比を取って同じ答えを得る.
    分数型は,隣り合う項の比として表すと隣接3項型になる.

    % G4
  \item \booknavref{節：等比型}{〔復習：等比型〕}\BookMathLayout{}{\\}
    第2項から第5項までに$r$を三回掛けるので,$48=-6r^3$である.
    $r$は実数だから$r=-2$, $a_1=\frac{a_2}{r}=3$となり,
    \[
      a_n=3(-2)^{n-1}.
    \]

    % F4
  \item \booknavref{節：階差型}{〔復習：階差型〕}\BookMathLayout{}{\\}
    因数の差が$2$なので,\bookterm{partial}{部分分数分解}すると,
    \[
      \frac2{(k+1)(k+3)}=\frac1{k+1}-\frac1{k+3}.
    \]
    \[
      \begin{aligned}
        a_n&=\sum_{k=1}^{n-1}\left(\frac1{k+1}-\frac1{k+3}\right)\\
        &=\frac12+\frac13-\frac1{n+1}-\frac1{n+2}.
      \end{aligned}
    \]
    二つの和が二項ずつずれるため,両端に二項ずつ残る.最終式は$n=1,2$でも直接確かめられる.

    % T4
  \item \booknavref{節：隣接3項型}{〔復習：隣接3項型〕}\BookMathLayout{}{\\}
    定数解$\alpha=2\alpha+3\alpha-8$から$2$を引く.
    $b_n=a_n-2$は$b_{n+2}=2b_{n+1}+3b_n$, $b_1=1,b_2=3$を満たす.
    \bookterm{characteristic}{特性方程式}の根$3,-1$のうち$3$を使い,$c_n=b_{n+1}-3b_n$とおくと,
    \[
      c_{n+1}=-c_n,\quad c_1=0.
    \]
    よって$c_n=0$から$b_{n+1}=3b_n$となり,$a_n=2+3^{n-1}$を得る.
    一方の\bookterm{auxiliary}{補助数列}が0になるので,二本目を解く前に復元できる.
    \par\smallskip\noindent\textbf{【別解】}\par
    2項間の漸化式$a_{n+1}=3a_n-4$が成り立つことを帰納法で示す.
    $a_2=3a_1-4$であり,この関係が$k$で成り立つなら,
    \[
      a_{k+2}-(3a_{k+1}-4)=-a_{k+1}+3a_k-4=0.
    \]
    あとは$a_{n+1}-2=3(a_n-2)$を解けばよい.

    % U4
  \item \booknavref{節：連立型}{〔復習：連立型〕}\BookMathLayout{}{\\}
    $s_n=a_n+b_n$, $d_n=a_n-b_n$とおく.
    和では定数項が消え,差では$2$となるので,
    \BookMathLayout{\[
        s_{n+1}=3s_n,\quad d_{n+1}=d_n+2,
        \quad s_1=1,\ d_1=-1.
      \]}{\[
        \begin{aligned}
          s_{n+1}&=3s_n,\quad d_{n+1}=d_n+2,\\
          s_1&=1,\quad d_1=-1.
        \end{aligned}
      \]}
    よって$s_n=3^{n-1}$, $d_n=2n-3$であり,
    \[
      a_n=\frac{s_n+d_n}2,\qquad b_n=\frac{s_n-d_n}2
    \]
    から解答の式を得る.

    % L4
  \item \booknavref{節：対数型}{〔復習：対数型〕}\BookMathLayout{}{\\}
    $a_1>0$と正の項から次も正になることより,全項が正である.
    $b_n=\log_2a_n$とおくと$b_{n+1}=1+\frac{b_n}{2}$, $b_1=0$となる.
    \bookterm{fixed}{不動点}$2$を引けば,
    \[
      b_n-2=-2\left(\frac12\right)^{n-1},
      \qquad a_n=2^{\,2-2^{2-n}}.
    \]
    \par\smallskip\noindent\textbf{【別解：\bookterm{fixed}{不動点}で割る】}\par
    正の\bookterm{fixed}{不動点}$4$で割り,$c_n=\frac{a_n}{4}$とおくと$c_{n+1}=\sqrt{c_n}$, $c_1=\frac{1}{4}$である.
    指数が毎回半分になるので,
    \[
      c_n=\left(\frac14\right)^{(\frac12)^{n-1}},
      \qquad a_n=4c_n.
    \]

    % C4
  \item \booknavref{節：特性方程式型}{〔復習：特性方程式型〕}\BookMathLayout{}{\\}
    \bookterm{fixed}{不動点}$3$との差$b_n=a_n-3$は$b_{n+1}=4b_n$, $b_1=t-3$を満たす.
    したがって$a_n=3+(t-3)4^{n-1}$である.
    $t-3$で割っていないため$t=3$も含み,このときは定\bookterm{sequence}{数列}となる.
    逆に定\bookterm{sequence}{数列}には$a_2=a_1$,すなわち$4t-9=t$が必要なので,$t=3$だけである.

    % R4
  \item \booknavref{節：階比型}{〔復習：階比型〕}\BookMathLayout{}{\\}
    初項も係数も正なので全項が正である.係数を順に掛けると,
    \[
      \begin{aligned}
        a_n&=2\prod_{k=1}^{n-1}2^{2k+1}\\
        &=2^{\,1+\sum_{k=1}^{n-1}(2k+1)}=2^{n^2}.
      \end{aligned}
    \]
    積を指数の和へ直した.空積なら初項も含む.
    \par\smallskip\noindent\textbf{【別解1】}\par
    $b_n=\log_2a_n$なら$b_{n+1}=b_n+2n+1$, $b_1=1$である.
    \[
      b_n=1+\sum_{k=1}^{n-1}(2k+1)=n^2
    \]
    と求め,$a_n=2^{b_n}$で戻す.
    \par\smallskip\noindent\textbf{【別解2】}\par
    $2n+1=(n+1)^2-n^2$に合わせれば,
    \[
      \frac{a_{n+1}}{2^{(n+1)^2}}=\frac{a_n}{2^{n^2}}=1
    \]
    となる.最後の値は初項から決まる.
    積と\bookterm{logarithm}{対数}は同じ和を計算し,別解2はその和の計算も省く.
    初項が負でも積と別解2は使えるが,$\log_2a_n$は使えない.
    \sectionref*{節：操作を選ぶ練習}の第3問に対応する.

    % D4
  \item \booknavref{節：等差型}{〔復習：等差型〕}\BookMathLayout{}{\\}
    第3項から第8項までに公差を五回加えるので,$-7-8=5d$より$d=-3$である.
    $a_3=a_1+2d$から$a_1=14$となり,$a_n=14-3(n-1)=17-3n$を得る.

    % M4
  \item \booknavref{節：階差等比複合型}{〔復習：階差等比複合型〕}\BookMathLayout{}{\\}
    係数と付加項の指数の底がともに$2$なので,$b_n=\frac{a_n}{2^{n-1}}$とおく.
    $b_{n+1}=b_n+3$, $b_1=2$から$b_n=3n-1$であり,
    \[
      a_n=(3n-1)2^{n-1}.
    \]
    底が等しいため,割った残りは\bookterm{fixed}{不動点}を使わずに等差型で解ける.
    \par\smallskip\noindent\textbf{【別解：形を置く】}\par
    付加項$3\cdot2^n$と同じ形$A\cdot2^{n-1}$は付加項のない式の解なので,$n$を掛けて$g(n)=An\cdot2^{n-1}$と置く.
    $g(n+1)=2g(n)+3\cdot2^n$に代入すると$A=3$であり,
    $a_n=C\cdot2^{n-1}+3n\cdot2^{n-1}$に初項を入れれば$C=-1$となる.
    \sectionref*{節：答えの形を見通す}の考え方である.

    % U5
  \item \booknavref{節：連立型}{〔復習：連立型〕}\BookMathLayout{}{\\}
    \BookMathLayout{第1式を進めて第2式を代入すると,$a_{n+2}=a_{n+1}+2a_n$となる.}{第1式を進めて第2式を代入すると,\newline
      $a_{n+2}=a_{n+1}+2a_n$となる.}
    \bookterm{characteristic}{特性方程式}の根は$2,-1$であり,
    $a_n=C2^{n-1}+D(-1)^{n-1}$とおく.
    $a_1=a_2=1$より$C=\frac{2}{3},D=\frac{1}{3}$だから,
    \[
      a_n=\frac{2^n+(-1)^{n-1}}3.
    \]
    $n\ge2$では$b_n=a_{n-1}$なので,
    \[
      b_n=\frac{2^{n-1}-(-1)^{n-1}}3.
    \]
    この最終式は$n=1$でも$b_1=0$に合う.
    \par\smallskip\noindent\textbf{【別解】}\par
    $c_n=a_n+b_n$, $d_n=a_n-2b_n$とおくと,
    \[
      c_{n+1}=2c_n,\quad d_{n+1}=-d_n,
      \quad c_1=d_1=1.
    \]
    よって$c_n=2^{n-1}$, $d_n=(-1)^{n-1}$であり,
    $3b_n=c_n-d_n$, $3a_n=2c_n+d_n$で直接戻せる.

    % R5
  \item \booknavref{節：階比型}{〔復習：階比型〕}\BookMathLayout{}{\\}
    添字を一つ進めると,$(n-1)!$には$n$が,$n!$には$n+1$が掛かる.
    2つの階乗を組み合わせれば,$n$によって変わる係数$n(n+1)$を作れる.
    $n(n+1)$を約分するために両辺を$n!(n+1)!$で割り,\BookMathLayout{}{\\}%
    $b_n=\frac{a_n}{(n-1)!n!}$とおくと,
    \[
      b_{n+1}=\frac{n(n+1)a_n}{2n!(n+1)!}
      =\frac12b_n.
    \]
    $0!=1$より$b_1=1$なので,\BookMathLayout{}{\\}%
    $a_n=\frac{(n-1)!n!}{2^{n-1}}$を得る.
    \par\smallskip\noindent\textbf{【別解】}\par
    係数を分けて掛ければ,
    \[
      \begin{aligned}
        a_n&=\frac1{2^{n-1}}\prod_{k=1}^{n-1}k
        \prod_{k=1}^{n-1}(k+1)\\
        &=\frac{(n-1)!n!}{2^{n-1}}.
      \end{aligned}
    \]
    $n=1$では二つの空積を$1$とする.

    % Q5
  \item \booknavref{節：一次分数型}{〔復習：一次分数型〕}\BookMathLayout{}{\\}
    \bookterm{fixed}{不動点}の方程式は$(x-1)^2=0$である.
    二つの異なる差の比は作れないので,
    \[
      a_{n+1}-1=\frac{a_n-1}{a_n}
    \]
    を調べる.$a_1>1$から全項で$a_n>1$となり,元の分母と差はともに0でない.
    $b_n=\frac{1}{a_n-1}$とおくと,
    \[
      b_{n+1}=b_n+1,\quad b_1=1,
      \quad a_n=1+\frac1n=\frac{n+1}n.
    \]
    重なる\bookterm{fixed}{不動点}からの差の逆数が等差型になった.
    \par\smallskip\noindent\textbf{【別解】}\par
    最初の三項$\frac{2}{1},\frac{3}{2},\frac{4}{3}$から同じ式を予想する.
    初項に合い,$a_k=\frac{k+1}{k}$なら$a_{k+1}=2-\frac{k}{k+1}=\frac{k+2}{k+1}$となるので,
    \bookterm{induction}{数学的帰納法}で示せる.
    \par\smallskip\noindent\textbf{【初項を変えたとき】}\par\nopagebreak[4]
    初項$1$なら定数数列となり差の逆数は使えない.
    初項$0$なら第2項が定義できず,$\frac{1}{2}$なら次項が$0$になって第3項が定義できない.
    置き換えに必要な条件と,元の漸化式に必要な条件を区別する.

    % G5
  \item \booknavref{節：等比型}{〔復習：等比型〕}\BookMathLayout{}{\\}
    $a_1=0$であり,$0$からは次も$0$となるので全項が$0$である.
    \bookterm{geometric}{等比数列}の式でも$a_n=0\cdot7^{n-1}=0$となるが,
    項で割って$\frac{a_{n+1}}{a_n}$を作ってはいけない.

    % L5
  \item \booknavref{節：対数型}{〔復習：対数型〕}\BookMathLayout{}{\\}
    正の項から次も正になるので全項の\bookterm{logarithm}{対数}を取れる.
    $b_n=\log_2a_n$とおくと,
    \[
      b_{n+1}=2b_n-n+1,\quad b_1=2.
    \]
    $b_{n+1}=2b_n-n+1$を満たす一次式$g(n)=un+v$を探すと,$u=2u-1$, $u+v=2v+1$より$g(n)=n$である.
    $c_n=b_n-n$は$c_{n+1}=2c_n$, $c_1=1$を満たすので,
    \[
      b_n=n+2^{n-1},\qquad a_n=2^{\,n+2^{n-1}}.
    \]
    \par\smallskip\noindent\textbf{【別解：$2^n$で割る】}\par
    \bookterm{logarithm}{対数}で$n$を引く操作は,元の項を$2^n$で割ることに対応する.
    $d_n=\frac{a_n}{2^n}$とおけば,
    \[
      d_{n+1}=\frac{a_n^2}{2^{2n}}=d_n^2,\quad d_1=2.
    \]
    よって$d_n=2^{2^{n-1}}$であり,$a_n=2^nd_n$で同じ答えを得る.

    % F5
  \item \booknavref{節：階差型}{〔復習：階差型〕}\BookMathLayout{}{\\}
    \bookterm{difference}{階差}を足し,公比$-1$の和を使うと,
    \[
      a_n=1+\sum_{k=1}^{n-1}(-1)^k
      =\frac{1+(-1)^{n-1}}2.
    \]
    空の和なら初項も含み,奇数番目は$1$,偶数番目は$0$となる.
    \par\smallskip\noindent\textbf{【別解】}\par
    二回分をまとめると$a_{n+2}=a_n+(-1)^n+(-1)^{n+1}=a_n$である.
    $a_1=1,a_2=0$なので$1,0$を繰り返し,同じ式で表せる.

    % C5
  \item \booknavref{節：特性方程式型}{〔復習：特性方程式型〕}\BookMathLayout{}{\\}
    $n=1,2$の漸化式から$5=8p+q$, $\frac{7}{2}=5p+q$となる.
    引いて$q$を消すと$p=\frac{1}{2}$, $q=1$である.
    \bookterm{fixed}{不動点}$2$を引けば,
    \[
      a_{n+1}-2=\frac12(a_n-2),\quad a_1-2=6
    \]
    より$a_n=2+6\left(\frac{1}{2}\right)^{n-1}$となる.
    \par\smallskip\noindent\textbf{【別解：\bookterm{difference}{階差}から\bookterm{general}{一般項}へ】}\par
    添字を1つずらした式と引くと,$d_n=a_{n+1}-a_n$は$d_{n+1}=\frac{d_n}{2}$, $d_1=-3$を満たす.
    よって,
    \[
      a_n=8-3\sum_{k=1}^{n-1}\left(\frac12\right)^{k-1}
      =2+6\left(\frac12\right)^{n-1}.
    \]
    空の和なら初項も含む.

    % T5
  \item \booknavref{節：隣接3項型}{〔復習：隣接3項型〕}\BookMathLayout{}{\\}
    定数を入れると左辺は$\alpha-3\alpha+2\alpha=0$となるため,定数解はない.
    $d_n=a_{n+1}-a_n$でまとめると,
    \[
      d_{n+1}=2d_n+1,\quad d_1=1.
    \]
    $d_{n+1}+1=2(d_n+1)$から$d_n=2^n-1$であり,
    \[
      a_n=1+\sum_{k=1}^{n-1}(2^k-1)=2^n-n.
    \]
    空の和なら初項も含む.\bookterm{difference}{階差}を取り,さらに$1$を加えて等比型にした.
    \par\smallskip\noindent\textbf{【別解】}\par
    $-n$は$-(n+2)-3\{-(n+1)\}+2(-n)=1$を満たす.
    そこで$b_n=a_n+n$とおくと,
    \[
      b_{n+2}-3b_{n+1}+2b_n=0,\quad b_1=2,\ b_2=4.
    \]
    $c_n=b_{n+1}-2b_n$は$c_{n+1}=c_n$, $c_1=0$を満たすので$b_n=2^n$となる.
    定数の代わりに一次式を引いても,右辺を消せる.
    これは\sectionref*{節：答えの形を見通す}で,定数の候補に$n$を掛けた形を使う例である.

    % D5
  \item \booknavref{節：等差型}{〔復習：等差型〕}\BookMathLayout{}{\\}
    両辺を$4$で割ると$a_{n+1}=a_n$であり,全項が初項$\frac{7}{4}$に等しい.
    公差$0$の\bookterm{arithmetic}{等差数列}としても$a_n=\frac{7}{4}+(n-1)\cdot0=\frac{7}{4}$となる.

    % M5
  \item \booknavref{節：階差等比複合型}{〔復習：階差等比複合型〕}\BookMathLayout{}{\\}
    $(-1)^{n-1}\ne0$なので$b_n=\frac{a_n}{(-1)^{n-1}}$とおく.
    $b_{n+1}=-2b_n+3$, $b_1=0$から\bookterm{fixed}{不動点}$1$を引けば,
    \[
      b_n=1-(-2)^{n-1}.
    \]
    $(-1)^{n-1}(-2)^{n-1}=2^{n-1}$を使って戻すと,
    \[
      a_n=(-1)^{n-1}-2^{n-1}.
    \]
    \par\smallskip\noindent\textbf{【別解：交代する項を引く】}\par
    $g_n=(-1)^{n-1}$も$g_{n+1}=2g_n+3(-1)^n$を満たす.
    与式から引くと$a_{n+1}-g_{n+1}=2(a_n-g_n)$であり,
    $a_1-g_1=-1$より$a_n-g_n=-2^{n-1}$を得る.

    % S5
  \item \booknavref{節：総和型}{〔復習：総和型〕}\BookMathLayout{}{\\}
    隣り合う総和の式を引くと,
    \[
      a_{n+1}=\frac{n+3}3a_{n+1}-\frac{n+2}3a_n
    \]
    より$na_{n+1}=(n+2)a_n$である.係数を順に掛ければ,
    \[
      a_n=2\prod_{k=1}^{n-1}\frac{k+2}{k}=n(n+1).
    \]
    $n=1$の空積は$1$とする.
    \par\smallskip\noindent\textbf{【別解】}\par
    同じ漸化式を
    \[
      \frac{a_{n+1}}{(n+1)(n+2)}=\frac{a_n}{n(n+1)}
    \]
    と書けば,この比は初項での値$1$に等しい.
\end{enumerate}
